\section{Methodology}
\label{sec:method}

\Cref{fig:pipeline} shows the pipeline of this study, marking each element as inherited, adapted for compositional
decision variables, or new here.

\begin{figure}[t]
\centering
\includegraphics[width=\textwidth]{fig01_lineage_pipeline.pdf}
\caption{The predecessor DoE--RSM--NBI framework (top) and its transfer to classifier ensemble weighting (bottom).
Blue: inherited. Orange: adapted because the decision variables are compositional and the responses are
machine-learning metrics. Green: introduced in this study.}
\label{fig:pipeline}
\end{figure}

\subsection{The ensemble-weighting problem is a mixture problem}
\label{sec:method:problem}

Let $M$ probabilistic base classifiers produce, for a sample $x$, probability estimates
$p_1(x),\dots,p_M(x)$ in $[0,1]$. A weighted soft-voting ensemble predicts
\begin{equation}
p_{\wvec}(x) \;=\; \sum_{i=1}^{M} w_i\, p_i(x) \;=\; \mathbf{P}\wvec ,
\label{eq:ensemble}
\end{equation}
where $\mathbf{P} \in [0,1]^{n \times M}$ collects the base-model probabilities for $n$ samples. For
$p_{\wvec}(x)$ to remain a probability for every input, the weights must satisfy
\begin{equation}
w_i \ge 0 \quad (i = 1,\dots,M), \qquad \sum_{i=1}^{M} w_i = 1 ,
\label{eq:simplex}
\end{equation}
that is, $\wvec \in \simplex$, the unit simplex in $\mathbb{R}^{M}$ and the canonical constraint set of a
\emph{mixture experiment}: the components are the classifiers, the vertices are single models used alone, and the
centroid is uniform soft voting. The sum-to-one constraint makes an intercept model singular, so the formalism is the
mixture design with a Scheff\'e canonical polynomial, not a factorial design. This correspondence, rather than any
performance argument, is why the framework transfers at all.

Three objectives are considered, all functions of $\wvec$ through \cref{eq:ensemble}:
\begin{align}
f_1(\wvec) &= -\operatorname{AUC}\!\left(y, \mathbf{P}\wvec\right) && \text{(discrimination, maximized)},\\
f_2(\wvec) &= \frac{1}{n}\sum_{k=1}^{n}\Big[-y_k \log p_{\wvec}(x_k) - (1-y_k)\log\!\big(1-p_{\wvec}(x_k)\big)\Big] && \text{(calibration, minimized)},\\
f_3(\wvec) &= c(\wvec) && \text{(inference cost, minimized)}.
\end{align}
The log-loss $f_2$ is convex and smooth in $\wvec$, whereas the area under the ROC curve is a rank statistic: $f_1$ is
piecewise constant in $\wvec$.

\subsection{Two cost models: a relaxation and a deployment cost}
\label{sec:method:cost}

Let $c_i$ be the measured inference cost of model $i$ (\cref{sec:protocol:models}). Deployment cost is a step
function of the weight vector: a model receiving \emph{any} non-negligible weight must be loaded and executed; one
receiving none costs nothing. This gives the \emph{support cost}
\begin{equation}
c_{\mathrm{sup}}(\wvec) \;=\; \sum_{i=1}^{M} c_i \,\mathbf{1}\!\left[w_i > \varepsilon\right], \qquad \varepsilon = 10^{-3},
\label{eq:supportcost}
\end{equation}
which is discontinuous; every optimizer here instead sees the continuous relaxation
\begin{equation}
c_{\mathrm{w}}(\wvec) \;=\; \sum_{i=1}^{M} w_i c_i ,
\label{eq:weightedcost}
\end{equation}
the natural linear mixture response, which coincides with $c_{\mathrm{sup}}$ at the vertices but not on the
interior. We therefore optimize \cref{eq:weightedcost} and score the \emph{same} solutions under
\cref{eq:supportcost}, reporting every comparison under both; \cref{sec:results:cost} shows that the choice changes
which method wins.

\subsection{Mixture design and Scheff\'e surrogates}
\label{sec:method:scheffe}

\paragraph{Design.} With $M = 5$ the design is a fixed 66-run mixture design: the union of the $\{5,3\}$ and
$\{5,2\}$ simplex lattices, the overall centroid, the 5 axial points, the 5 quaternary centroids and the 10 midpoints
between the overall centroid and the ternary centroids. It therefore places design points on the 5 pure vertices, all
10 binary edges, the ternary and quaternary faces and uniform voting, supporting the surrogate on every face a sparse
ensemble could occupy.

\paragraph{Models.} For a response $y$ the Scheff\'e canonical polynomials of order $1$, $2$ and the special cubic are
\begin{align}
\text{linear:}\quad & \hat y = \sum_i \beta_i w_i, \\
\text{quadratic:}\quad & \hat y = \sum_i \beta_i w_i + \sum_{i<j} \beta_{ij} w_i w_j, \label{eq:scheffequad}\\
\text{special cubic:}\quad & \hat y = \sum_i \beta_i w_i + \sum_{i<j} \beta_{ij} w_i w_j + \sum_{i<j<k} \beta_{ijk} w_i w_j w_k,
\end{align}
fitted by ordinary least squares without an intercept and without pure quadratic terms, which $\sum_i w_i = 1$ makes
redundant. No term elimination is performed: the canonical form is kept intact so that $\beta_i$ is the predicted
response at vertex $i$.

\paragraph{External validation and the reliability gate.} Each replication also evaluates 100 compositions drawn
from $\mathrm{Dirichlet}(\mathbf{1})$ (60) and $\mathrm{Dirichlet}(0.5\cdot\mathbf{1})$ (40) that are not in the design and never used for fitting, giving an
external $\rext$ and Spearman $\rho_{\mathrm{ext}}$. \emph{Order selection} uses a parsimony rule --- the lowest
order whose external RMSE is within 10\% of the best --- so it is driven by out-of-design accuracy rather than
in-sample fit. The \emph{same} 100 compositions serve both steps: the parsimony rule picks the order on them and the
gate then scores the selected order on them. The gate is therefore not independent of order selection --- the order it
judges was chosen to do well on exactly these points --- and its pass rate should be read as an optimistic estimate of
out-of-design adequacy. We report alongside it the external RMSE of the selected surface in the response's own units
and the observed range of that response over the same 100 points (\cref{tab:gate}), because on a nearly flat surface
$\rext$ is ill-posed --- its denominator is the variance of an almost constant response --- so the RMSE read against
the range is the more informative adequacy measure. \emph{Admissibility} uses a \emph{reliability gate} whose thresholds were fixed and committed to
version control before the replicated benchmark was run, and never adjusted afterwards: a surface is reliable only if
\begin{equation}
\rext \ge 0.5 \quad\text{and}\quad \rho_{\mathrm{ext}} \ge 0.9 .
\label{eq:gate}
\end{equation}
We describe this as pre-specification rather than preregistration: there is no
external registry entry, and the evidence is the version-controlled commit history. The gate is a hypothesis under
test, not a filter applied to the results: surrogate-based optimization is run and reported \emph{regardless} of the
gate outcome, so that its predictive value can be measured (\cref{sec:results:gate}).

\subsection{Normal Boundary Intersection on the simplex}
\label{sec:method:nbi}

NBI \citep{das1998nbi} was introduced to obtain a distributed representation of the Pareto boundary
(\cref{sec:rel:lineage}). Whether that distribution survives the map into revalidated objective space is one of the
questions this paper tests, and \cref{sec:results:metamodelfree} reports that in our setting it does not. For $m$
objectives $F = (F_1,\dots,F_m)$ let $\wvec^{*(k)}$ minimize $F_k$; the payoff matrix $\Phi$ has columns $F(\wvec^{*(k)})$, its diagonal is the utopia
point $F^{\mathrm{u}}$, and objectives are normalized by the utopia and pseudo-nadir points extracted from $\Phi$.
Writing $\bar\Phi$ for the normalized payoff matrix and $\hat{n} = -\bar\Phi\mathbf{1}$ for the quasi-normal direction,
each $\beta$ on the unit simplex of \emph{scalarization} weights defines the subproblem
\begin{equation}
\max_{\wvec,\,t} \; t \quad \text{s.t.} \quad \bar\Phi\beta + t\,\hat{n} = \bar F(\wvec), \quad \wvec \in \simplex .
\label{eq:nbi}
\end{equation}
The simplex constraint is eliminated by optimizing over $M-1$ free variables with
$w_M = 1 - \sum_{i<M} w_i$, with a deterministic feasibility-recovery step for iterates that leave the simplex: the
free variables are clipped to $[0,1]$ and, if they still sum above one, renormalized onto the face $w_M = 0$. This is
a documented repair, not a Euclidean projection onto the simplex, and a candidate recovered this way is flagged
infeasible and excluded from the certified set. Without it, subproblems near a vertex terminate outside the feasible
region. For vertex values of $\beta$ the solution is the corresponding anchor, returned directly rather than
re-derived. All variants solve 66 subproblems per replication on the same $\beta$ lattice, so the arms are
compared at identical points of the CHIM.

\paragraph{The three variants.} The variants are defined by two methodological factors --- the source of the
objectives and the source of the anchors:

\begin{center}
\footnotesize
\setlength{\tabcolsep}{4pt}
\begin{tabular}{llll}
\toprule
Variant & Objectives in \cref{eq:nbi} & Anchors / payoff matrix & Real evaluations, standalone \\
\midrule
NBI-A & Scheff\'e surfaces & optima of the Scheff\'e surfaces & $1.7\times10^{2}$ \\
NBI-B & Scheff\'e surfaces & real out-of-fold optima & $\approx 4.1\times10^{4}$ \\
NBI-C & real out-of-fold objectives & real out-of-fold optima & $\approx 4.6\times10^{5}$ \\
\bottomrule
\end{tabular}
\end{center}

NBI-A is the lineage construction: everything, including the payoff matrix, comes from the metamodel. NBI-B changes
only the anchors, which become the real optima of \cref{sec:method:refs}; NBI-C removes the metamodel entirely.
The last column charges every variant for each real objective evaluation it needs \emph{standalone}: the 66 design
points and the 100 gate compositions for NBI-A; those plus the real anchors for NBI-B, dominated by the direct
ROC-AUC search over $4\times10^{4}$ sampled compositions (\cref{sec:method:refs}); and the real anchors plus its own
solve for NBI-C. Real anchors are not free. NBI-B consumes about $4\times10^{4}$ real evaluations that NBI-A does
not --- roughly a tenth of NBI-C's entire budget --- and its standalone wall clock is $1.2$--$12.7$ times NBI-A's,
which is why \cref{tab:compute} prices each arm for the stages it depends on rather than for its solver stage alone.
Because $f_1$ is piecewise constant, NBI-C uses finite-difference steps of $10^{-3}$, multistart on the simplex, and
accepts an iterate whose equality residual is below $10^{-3}$ even without certified optimality,
whereas NBI-A and NBI-B optimize smooth polynomials and require full SLSQP certification. This asymmetry follows from
the objective, not from tuning, and it bounds the causal reading. A versus B is not a single-factor contrast:
replacing the anchors simultaneously moves the utopia point, rescales both normalization axes through the payoff
matrix, reorients the CHIM and every quasi-normal direction, and---because the three vertex values of $\beta$ return
their anchor directly---substitutes three returned candidates. \Cref{sec:results:anchors} separates the last of these
from the rest. B versus C tests removal of
\emph{surrogate error in the interior} once real anchors are supplied, but is not a perfectly identical numerical
experiment, because changing the objective from a smooth polynomial to a piecewise-constant rank statistic
necessarily changes solver semantics. Solver outcomes are therefore reported separately for each arm.

\subsection{References, comparators and revalidation}
\label{sec:method:refs}

\paragraph{Single-objective references.} Each replication computes the best single model, uniform voting, logistic
stacking, the SLSQP log-loss optimum (convex, hence global), a derivative-free ROC-AUC optimum over $4\times10^{4}$
sampled compositions with local refinement, and the optima of the two Scheff\'e surfaces. The log-loss and AUC optima
double as the real anchors of NBI-B and NBI-C. SLSQP on the convex log-loss remains the correct tool for the
single-objective problem; NBI is not proposed as a replacement for it.

\paragraph{Comparators.} Two cheap alternatives are run on every replication: random weighted scalarization, which
draws 66 scalarization weight vectors from $\mathrm{Dirichlet}(\mathbf{1})$ and minimizes the corresponding weighted
sum of the \emph{same} Scheff\'e surfaces with the same real anchors used by NBI-B; and a 66-point random
$\mathrm{Dirichlet}(\mathbf{1})$ sample of the weight simplex. The second is matched to the NBI sets only in the
\emph{number of returned candidates} and is not compute-matched; it functions as a floor check on the geometry of the
reference front, not as an efficiency baseline.

\paragraph{Revalidation.} Every candidate from every method --- the three NBI arms, both comparators, the 66 design
runs and the references --- is re-evaluated on the exact out-of-fold objectives before any indicator is computed, with
non-dominated filtering applied to those revalidated values. No result in this paper is read off a surrogate.

\subsection{Empirical Pareto reference and quality indicators}
\label{sec:method:reference}

The true Pareto set of \cref{eq:ensemble} over $\simplex$ is unknown, so each replication constructs an
\emph{empirical Pareto reference}: at least $10^{5}$ compositions
from $\mathrm{Dirichlet}(\mathbf{1})$ and $\mathrm{Dirichlet}(0.3\cdot\mathbf{1})$ in equal parts, augmented with the
simplex lattice, all vertices and edges, and a 40-cap $\varepsilon$-constraint sweep, evaluated on the real objectives
and filtered to the non-dominated set. An independent sample of $2\times10^{4}$ points checks convergence: if it
displaces more than 5\% of the front, the reference is enlarged and rechecked, up to three rounds. This sampled core
contains no output of the five methods compared in this study --- NBI-A, NBI-B, NBI-C, random weighted scalarization
and NSGA-II. It is not, however, free of optimization output: it carries the 40 $\varepsilon$-constraint solutions and
the seven single-objective references of \cref{sec:method:refs}, among them the real ROC-AUC and log-loss optima that
NBI-B and NBI-C receive as anchors and the two Scheff\'e-surface optima that NBI-A receives. On Santander the
$\varepsilon$-constraint sweep supplies a median of 38 of the roughly 135 core-front points and the single-objective
references a median of 4 (\texttt{tables/reference\_diagnostics.csv}). The final reference used for scoring
is the non-dominated union of that core with every candidate set of the replication, making it at least as good as
any method scored against it.

We call this an empirical reference, never a true Pareto front: it is a sample-based upper bound, and because
candidate sets contribute to it, a method contributing many points is partly graded against itself.

Approximation sets are scored on min--max normalized objectives by generational distance, inverted generational
distance, $\igdp$ \citep{ishibuchi2015modified}, Schott spacing, the exact hypervolume with reference point $1.1$ on
each axis, expressed as a ratio to the reference hypervolume, coverage, and the joint non-dominated fraction.

\subsection{The normalization box}
\label{sec:method:normbox}

Every indicator is computed on objectives mapped into $[0,1]$ by one box, and the box is not a fixed property of the
problem. Let $\mathcal{F}$ be the empirical reference front of the replication --- the non-dominated union of the
sampled core with every candidate set --- under one of the two cost definitions. The box is the componentwise
minimum and maximum of $\mathcal{F}$:
\begin{equation}
\ell_k \;=\; \min_{F \in \mathcal{F}} F_k, \qquad u_k \;=\; \max_{F \in \mathcal{F}} F_k, \qquad
\tilde F_k \;=\; \frac{F_k - \ell_k}{u_k - \ell_k}, \qquad k = 1,2,3 .
\label{eq:normbox}
\end{equation}
Three consequences follow, and each constrains how the numbers in \cref{sec:results} may be read.

\emph{(i) Within a replication the box is common to every arm.} All candidate sets of a replication are scored under
the same $\ell$ and $u$, so the paired contrast between two arms is a difference of commensurable quantities and is
unaffected by the choice of box.

\emph{(ii) The box changes when a method is added to the reference,} because a new candidate set can extend
$\mathcal{F}$. Adding the NSGA-II arm moves every indicator level slightly: on the identical candidate sets, NBI-C's
median hypervolume ratio is $0.989$, $0.983$, $0.982$ and $0.976$ in \cref{sec:results:metamodelfree}, computed
without NSGA-II in the reference, and $0.990$, $0.981$, $0.981$ and $0.977$ in \cref{tab:nsga2}, computed with it.
Absolute indicator levels are therefore only interpretable together with the set of methods that built the reference.

\emph{(iii) The box differs between replications,} since each replication has its own reference front. Absolute
levels are comparable within a replication but not across partitions, and a mean of $\Delta\hv$ over 30 partitions is
a mean over 30 different scales. Paired within-replication differences are unaffected, which is why medians, win
counts and per-partition distributions are reported alongside means; \cref{sec:results:cost} withdraws one earlier
between-partition reading that this rescaling, not the optimizers, produced.
